\centerline{\bf \Huge Полупризнак равенства треугольников}

\begin{flushright}
        \scriptsize{Посмотри на обратную сторону!}
\end{flushright}

\problem{1} Докажите, что если в треугольнике биссектриса и медиана совпадают, то такой треугольник - равнобедренный

\problem{2} В четырехугольнике $ABCD$ известно, что прямые $AB$ и $CD$ пересекаются, сторона $AB=23$. $M$ - середина стороны $AD$. Оказалось, что $BM=CM$ и $\angle ABM = \angle DCM$. Найдите сторону $CD$.

%{\it Рассмотрите треугольники $ABM$ и $DCM$}

\problem{3} На биссектрисе угла $B$ неравнобедренного треугольника $ABC$ отметили точку $D$ такую, что $AD = DC$. Найдите углы  $\triangle ADC$, если известно, что $\angle ABC = 20^{\circ}$.

%{\it Подсказка: надо сначала доказать, что вариант с тем, что точка $D$ лежит внутри треугольника невозможен. После рассмотреть треугольники $ABD$ и $CBD$, откуда получить, что $\angle BAD + \angle BCD = 180^{\circ}$ и откуда будет следовать, что $\angle ADC = 160^{\circ}$}

\problem{4} В треугольнике $ABC$ сторона $AB$ больше стороны $BC$ и биссектриса $AL$ равна стороне $AC$. На биссектрисе $AL$ отметили точку $K$ так, что $CK=BL$. Докажите, что $\angle CKL = \angle ABC$.

\problem{5} В остроугольном треугольнике $ABC$ на сторонах $AC$ и $AB$ отметили точки $K$ и $L$ соответственно, причём прямая $KL$ параллельна $BC$ и $KL = KC$.  На стороне $BC$ выбрана точка $M$ так, что  $\angle KMB = \angle BAC$.  Докажите, что  $KM = AL$.

%{\it Постройте точку $D$ на стороне $BC$ такую, что $\angle KDM = \angle KMD$, после этого рассмотрите треугольники $ALK$ и $DKM$ }

\problem{6} Через середину $S$ отрезка $MN$, концы которого лежат на боковых сторонах равнобедренного треугольника, проведена прямая, параллельная основанию треугольника и пересекающая боковые стороны в точках $K$ и $L$. Докажите, что $MK = ML$.

%{\it На прямой $KL$ отложите точку $D$ так, что $\angle MDK = \angle MKD$ и рассмотрите треугольники $MDS$ и $NLS$}

\newpage

\hspace{3mm}

\centerline{\bf \Huge Прямоугольный треугольник}

\begin{flushright}
        \scriptsize{Ха, повёлся!}
\end{flushright}

\problem{1} В треугольнике $ABC$ известно, что $\angle A = 60^{\circ}$ и $AB=2AC$. Найдите углы треугольника $ABC$.


\problem{2} Биссектриса угла $A$ треугольника $ABC$, у которого $\angle C = 30^{\circ}$, перпендикулярна медиане, проведенной из вершины $B$. Найдите угол $B$. 

\problem{3} В четырёхугольнике $ABCD$ углы $ABD$ и $ACD$ прямые. Точки $M$ и $K$ --- середины его сторон $AD$ и $BC$ соответственно. Докажите, что прямая $MK$ перпендикулярна одной из сторон этого четырёхугольника.

\problem{4} В треугольнике $A B C$ известно, что $\angle A=15^{\circ}$ и $\angle B=30^{\circ}$. На стороне $A B$ отметили такую точку $E$, что $\angle A C E=90^{\circ}$. Найдите отрезок $A E$, если $B C=2$ см.

\problem{5} В прямоугольном треугольнике $ABC$ провели высоту $CH$ к гипотенузе. Известно, что $\angle A = 15^{\circ}$ и $CH = 1$. Найдите длину $AB$.

\problem{6} В прямоугольнике $ABCD$ известно, что сторона $AB$ в два раза меньше стороны $AD$. На стороне $BC$ выбрана точка $Q$ так, что $AQ=BC$. Чему равен угол $DQC$? 